\documentclass[11pt,a4paper]{article}

% Self-contained technical handoff for the RVE PROM--ANN/HPROM--ANN workflow.
\usepackage[margin=25mm]{geometry}
\usepackage[T1]{fontenc}
\usepackage[utf8]{inputenc}
\usepackage{lmodern}
\usepackage{amsmath,amssymb,mathtools,bm}
\usepackage{booktabs,array,longtable,enumitem}
\usepackage{xcolor}
\usepackage[hidelinks]{hyperref}
\usepackage[nameinlink,capitalise]{cleveref}

\newcommand{\R}{\mathbb{R}}
\newcommand{\transpose}{^{\mathsf T}}
\newcommand{\norm}[1]{\left\lVert #1 \right\rVert}
\newcommand{\frobnorm}[1]{\left\lVert #1 \right\rVert_{\mathrm F}}
\newcommand{\avg}[1]{\left\langle #1 \right\rangle}
\newcommand{\vect}[1]{\mathbf{#1}}
\newcommand{\mat}[1]{\mathbf{#1}}
\newcommand{\set}[1]{\mathcal{#1}}
\newcommand{\code}[1]{\url{#1}}

\title{\textbf{PROM--ANN, HPROM--ANN, and MAW--ECM}\\
\large Technical formulation and implementation handoff from the finite-deformation RVE workflow}
\author{Project technical note}
\date{\today}

\begin{document}
\maketitle

\begin{abstract}
This document is the implementation-level description of the RVE workflow in
\url{RVE_homogenization_NeoHookean_using_Kratos}.  It explains the
complete chain from macro-strain trajectories and full-order snapshots to the
structured PROM--ANN coordinates, the intrusive HPROM--ANN correction, the
direct D-HPROM--ANN variant, and state-dependent MAW--ECM hyper-reduction.
The notation follows the WCCM 2026 slides: the primary coordinates are
\(\vect q\), the primary and secondary bases are \(\mat V\) and
\(\overline{\mat V}\), and the ANN predicts \(\overline{\vect q}\).

The note is intentionally self-contained.  A future implementation for
\code{burgers1d-rom-workbench} should read this file together with the source
files listed in \cref{sec:source-map}, then port the algorithmic roles rather
than copying RVE-specific details such as finite-deformation lifting or stress
homogenization.
\end{abstract}

\tableofcontents
\newpage

\section{Purpose, scope, and current RVE configuration}
\label{sec:scope}

The objective is to approximate the nonlinear solution map of a finite-
deformation representative volume element (RVE), while preserving an
intrusive equilibrium correction when desired and evaluating both equilibrium
and homogenized quantities from a very small set of elements.

The current RVE instance has the following numerical configuration.  These
numbers document the present experiment; they are \emph{not} assumptions of
the method.

\begin{center}
\begin{tabular}{@{}ll@{}}
\toprule
Quantity & Current value \\
\midrule
Total displacement degrees of freedom & \(4244\) \\
Free displacement degrees of freedom & \(N_f=3924\) \\
Strong Dirichlet degrees of freedom & \(320\) \\
Training snapshots & \(N_s=109460\) \\
Truncated POD dimension & \(n_{\rm tot}=39\) \\
Primary and secondary dimensions & \(n=3\), \(\bar n=36\) \\
Full RVE elements & \(N_e=990\) \\
Final MAW residual support & \(|\set Z_{\rm res}|=10\) \\
Final MAW strain and stress supports & \(|\set Z_{\varepsilon}|=|\set Z_{\sigma}|=10\) \\
Union of the three final supports & \(29\) elements \\
\bottomrule
\end{tabular}
\end{center}

There are two online models.
\begin{description}[leftmargin=33mm,style=nextline]
  \item[HPROM--ANN] The affine map supplies an initial \(\vect q^0\), then
  a reduced Galerkin/Newton correction solves equilibrium on the
  ANN-defined nonlinear manifold.  MAW--ECM accelerates the residual and the
  homogenization.
  \item[D-HPROM--ANN] The direct approximation sets \(\vect q=\vect q^0\),
  decodes the state once, and applies the hyper-reduced homogenization.  It
  skips the residual evaluation and all reduced Newton iterations.
\end{description}

\section{Full-order RVE problem}
\label{sec:fom}

\subsection{Macro input and affine boundary lifting}

For each loading state, the physical macro input is either a deformation
gradient \(\mat F\) or, equivalently for the present workflow, its Green--
Lagrange strain
\begin{equation}
  \mat E
  = \frac{1}{2}\left(\mat F\transpose\mat F-\mat I\right),
  \qquad
  \vect\mu
  =
  \begin{bmatrix}
    E_{xx} & E_{yy} & G_{xy}
  \end{bmatrix}\transpose
  \in\R^3.
  \label{eq:macro-input}
\end{equation}
The code reconstructs the chosen two-dimensional branch of \(\mat F\) from
\(\mat E\).  Let \(\vect X\) denote a reference material point and
\(\vect X_c\) the RVE reference-domain center.  The imposed affine
displacement is
\begin{equation}
  \vect u_{\rm aff}(\vect X;\vect\mu)
  = \left(\mat F(\vect\mu)-\mat I\right)(\vect X-\vect X_c).
  \label{eq:affine-lifting}
\end{equation}
The total displacement is decomposed as
\begin{equation}
  \vect u(\vect\mu)
  = \vect u_{\rm aff}(\vect\mu)+\vect s(\vect\mu),
  \label{eq:state-decomposition}
\end{equation}
where the fluctuation \(\vect s\) is zero on the strongly prescribed
Dirichlet degrees of freedom.  All reduced bases and coordinates below act on
the \(N_f\) free entries of \(\vect s\).

At element level, let \(\vect r_e(\vect u;\vect\mu)\) and
\(\mat K_e(\vect u;\vect\mu)\) denote the element residual and its tangent.
The full internal equilibrium residual and tangent are assembled as
\begin{equation}
  \vect r(\vect u;\vect\mu)=\sum_{e=1}^{N_e}\vect r_e(\vect u;\vect\mu),
  \qquad
  \mat K(\vect u;\vect\mu)=\sum_{e=1}^{N_e}\mat K_e(\vect u;\vect\mu).
  \label{eq:full-assembly}
\end{equation}
The full-order model (FOM) solves \(\vect r=\vect 0\), with the sign
convention used by the underlying finite-element assembly.  All reduced
equations must retain that convention consistently.

\subsection{Homogenized strain and stress}

Let \(A_e\) be the reference measure of element \(e\), \(A_0=\sum_e A_e\)
the reference RVE measure, and \(\avg{\cdot}_e\) the element Gauss-point
average.  The code returns the same strain and stress measures that are
stored by the finite-element constitutive law.  Their full-order
homogenizations have the generic form
\begin{equation}
  \overline{\mat E}
  = \frac{1}{A_0}\sum_{e=1}^{N_e} A_e\avg{\mat E_e}_e,
  \qquad
  \overline{\mat S}
  = \frac{1}{A_0}\sum_{e=1}^{N_e} A_e\avg{\mat S_e}_e.
  \label{eq:full-homogenization}
\end{equation}
The symbol \(\mat S\) in \cref{eq:full-homogenization} is a stress tensor;
the snapshot matrix later in the document is written in bold as \(\mat S\)
only within the POD section.  Context and dimensions distinguish the two.

\section{Training trajectories and FOM data}
\label{sec:trajectories}

\subsection{Macro-strain domain}

The current training domain is a box in the three physical macro-strain
components.  With a prescribed positive and negative extent in each
direction, it can be written as
\begin{equation}
  \mathcal D_\mu
  = [E_{xx}^{-},E_{xx}^{+}]
    \times[E_{yy}^{-},E_{yy}^{+}]
    \times[G_{xy}^{-},G_{xy}^{+}].
  \label{eq:parameter-domain}
\end{equation}
For the WCCM data, \texttt{stage0\_training\_trajectory.py} uses
\(E_{\max}=2\) and relative bounds
\begin{equation}
  (1,0.05,1,0.05,0.05,0.05),
\end{equation}
giving the asymmetric box used in the trajectory figures.

\subsection{Serpentine paths}

The loading set is deliberately a collection of continuous paths rather than
an unordered cloud.  Five shear layers are used,
\begin{equation}
  G_{xy}\in\{0,+G_1,+G_2,-G_1,-G_2\},
  \label{eq:shear-layers}
\end{equation}
and each layer contains two scans starting from the origin:
an upper scan and a lower scan.  Within each fixed-shear plane, the trajectory
visits a sequence of \(E_{yy}\) rows and alternates the direction of its
\(E_{xx}\) sweep.  Consequently, the default construction contains ten
continuous serpentine trajectories.  This gives coverage of the macro domain
while avoiding unrelated jumps between FOM states.

For a path with waypoints \(\{\vect\mu_\ell\}_{\ell=0}^{L}\), the actual FOM
path is piecewise linear,
\begin{equation}
  \vect\mu(\lambda)
  =(1-\lambda)\vect\mu_\ell+\lambda\vect\mu_{\ell+1},
  \qquad \lambda\in[0,1].
  \label{eq:piecewise-linear-path}
\end{equation}
The number of increments on a segment is proportional to its scaled length.
In the present script the visualization and increment metric use
\begin{equation}
  \widehat{\vect\mu}
  = [E_{xx},E_{yy},G_{xy}/2]\transpose,
  \qquad
  n_\ell
  =\max\!\left(1,
  \left\lceil
  n_{\rm ref}\,
  \frac{\norm{\widehat{\vect\mu}_{\ell+1}-\widehat{\vect\mu}_{\ell}}}
       {a_{\rm ref}}
  \right\rceil\right).
  \label{eq:segment-steps}
\end{equation}
The scaling is an RVE trajectory-design choice, not part of PROM--ANN.
Its purpose is to keep nonlinear increments manageable and maintain a
comparable spatial resolution on every path.

\subsection{Snapshot collection}

The FOM is run on every trajectory.  At each increment \(j\), it stores
\begin{equation}
  \left\{\vect u^{(j)},\;\vect\mu^{(j)},\;
  \overline{\mat E}^{(j)},\;\overline{\mat S}^{(j)}\right\}.
  \label{eq:snapshot-tuple}
\end{equation}
The affine-free free-DOF snapshot is
\begin{equation}
  \vect s^{(j)}
  = \left.\left(\vect u^{(j)}-\vect u_{\rm aff}(\vect\mu^{(j)})\right)
    \right|_{\rm free}
  \in\R^{N_f}.
  \label{eq:fluctuation-snapshot}
\end{equation}
The exact vanishing of the Dirichlet entries of this quantity is checked in
Stage 2.  It is essential: otherwise the reduced coordinates would mix a
known boundary lifting with the unknown fluctuation.

\section{POD and the structured coordinate system}
\label{sec:coordinates}

\subsection{Truncated POD representation}

Collect the snapshots as columns:
\begin{equation}
  \mat S
  =\begin{bmatrix}
    \vect s^{(1)} & \cdots & \vect s^{(N_s)}
  \end{bmatrix}
  \in\R^{N_f\times N_s}.
  \label{eq:snapshot-matrix}
\end{equation}
The POD of \(\mat S\) gives the truncated orthonormal basis
\(\mat V_{\rm tot}\in\R^{N_f\times n_{\rm tot}}\).  In the present RVE,
\(n_{\rm tot}=39\) captures \(1-9.01\times10^{-11}\) of the snapshot
energy.  Define the standard POD coordinate matrix
\begin{equation}
  \mat Q_{\rm tot}
  = \mat V_{\rm tot}\transpose\mat S
  \in\R^{n_{\rm tot}\times N_s},
  \qquad
  \mat S\approx\mat V_{\rm tot}\mat Q_{\rm tot}.
  \label{eq:standard-pod-coordinates}
\end{equation}

The first three standard POD coordinates are excellent for minimizing a
linear reconstruction error.  They need not be the best three coordinates
for regression from \(\vect\mu\): a nonlinear loading manifold can appear
curved, coupled, or folded when projected onto those first modes.  The
following construction therefore uses \emph{all} retained POD coordinates to
identify a three-dimensional subspace aligned with macro loading.

\subsection{Linear alignment with macro strain}

Collect the same macro inputs as columns,
\begin{equation}
  \mat M
  = \begin{bmatrix}
    \vect\mu^{(1)} & \cdots & \vect\mu^{(N_s)}
  \end{bmatrix}
  \in\R^{3\times N_s}.
  \label{eq:macro-snapshot-matrix}
\end{equation}
The least-squares alignment matrix is
\begin{equation}
  \mat T
  = \underset{\widehat{\mat T}\in\R^{3\times n_{\rm tot}}}
            {\operatorname*{arg\,min}}
      \frobnorm{\widehat{\mat T}\mat Q_{\rm tot}-\mat M}^{2}.
  \label{eq:fit-T}
\end{equation}
Thus \(\mat T\mat Q_{\rm tot}\) is the best linear prediction of the three
macro-strain components available from the entire retained POD coordinate
set.  The implementation stores the row-sample equivalent of
\cref{eq:fit-T}: \(\mat Q_{\rm tot}\transpose\mat T\transpose\approx
\mat M\transpose\).

Map the aligned directions to physical free-DOF space and factor them by a
compact singular value decomposition,
\begin{equation}
  \mat V_{\rm tot}\mat T\transpose
  = \mat U\mat\Sigma\mat R\transpose
  = \mat V\mat A,
  \qquad
  \mat V=\mat U\in\R^{N_f\times3},
  \qquad
  \mat A=\mat\Sigma\mat R\transpose\in\R^{3\times3}.
  \label{eq:VA-factorization}
\end{equation}
The matrix \(\mat V\) is orthonormal.  The non-orthogonal scaling and
rotation required by the macro-aligned directions are retained in \(\mat A\).
This is why the state decoder contains \(\mat V\mat A\vect q\), not merely
\(\mat V\vect q\).

\subsection{Primary and secondary coordinates}

For an affine-free state \(\vect s\), define the structured primary
coordinates by
\begin{equation}
  \vect q
  = \mat A^{-1}\mat V\transpose\vect s\in\R^3.
  \label{eq:structured-q}
\end{equation}
For the whole data set this gives
\begin{equation}
  \mat V\transpose\mat S=\mat A\mat Q,
  \qquad
  \mat Q=\mat A^{-1}\mat V\transpose\mat S
  \in\R^{3\times N_s}.
  \label{eq:structured-Q}
\end{equation}
The remaining directions are an orthonormal complement in the retained POD
space.  In the code, if \(\mat C_m=\mat V_{\rm tot}\transpose\mat V\), then
\(\mat C_s\) spans the complement of \(\mat C_m\), and
\begin{equation}
  \overline{\mat V}=\mat V_{\rm tot}\mat C_s,
  \qquad
  \overline{\mat Q}=\overline{\mat V}\transpose\mat S
  \in\R^{36\times N_s}.
  \label{eq:secondary-coordinates}
\end{equation}
The truncated-POD reconstruction is therefore equivalently
\begin{equation}
  \mat S
  \approx
  \mat V\mat A\mat Q
  +\overline{\mat V}\overline{\mat Q}.
  \label{eq:split-reconstruction}
\end{equation}

\paragraph{Important distinction.}
The coordinates \(\vect q\) are \emph{not} the first three raw standard POD
coordinates.  They are the coordinates of a three-dimensional subspace
identified from all 39 retained modes through \cref{eq:fit-T} and then
factored through \cref{eq:VA-factorization}.  This is the precise meaning of
``structured coordinates better suited to regression.''

\section{PROM--ANN state approximation}
\label{sec:prom-ann}

\subsection{Affine primary-coordinate initialization}

The map from macro strain to the structured primary coordinate is very nearly
affine in the current data.  Define the row-sample design matrix
\begin{equation}
  \mat M_{\rm aff}
  =\begin{bmatrix}\mat M\transpose&\vect 1\end{bmatrix}
  \in\R^{N_s\times4}.
  \label{eq:maff}
\end{equation}
The coefficient matrix \(\mat B_{\rm aff}\in\R^{4\times3}\) is obtained by
\begin{equation}
  \mat B_{\rm aff}
  =\underset{\mat B\in\R^{4\times3}}
           {\operatorname*{arg\,min}}
   \frobnorm{\mat M_{\rm aff}\mat B-\mat Q\transpose}^{2}.
  \label{eq:Baff}
\end{equation}
For one macro state, the affine prediction is
\begin{equation}
  \left(\vect q^0(\vect\mu)\right)\transpose
  =\begin{bmatrix}\vect\mu\transpose&1\end{bmatrix}\mat B_{\rm aff}.
  \label{eq:q0}
\end{equation}
In the current RVE, the relative fit error of this map is
\(6.88\times10^{-4}\).  This makes \(\vect q^0\) an excellent initialization
for an intrusive solve and a meaningful direct prediction.

The role of \(\mat A\) deserves a precise statement.  It influences
\(\vect q\) through \cref{eq:structured-q}, so it influences the data used to
fit \(\mat B_{\rm aff}\).  However, \(\mat A\) does \emph{not} appear
explicitly in the least-squares problem \cref{eq:Baff}; that problem directly
fits \(\vect\mu\mapsto\vect q\).

\subsection{ANN closure of secondary coordinates}

The closure network learns the discarded coordinates as a function of the
three structured coordinates,
\begin{equation}
  \overline{\vect q}
  \approx\mathcal N_\theta(\vect q),
  \qquad
  \vect q\in\R^3,
  \qquad
  \overline{\vect q}\in\R^{36}.
  \label{eq:ann-closure}
\end{equation}
It is trained from the pairs
\(\{\vect q^{(j)},\overline{\vect q}^{(j)}\}\).  The production solver uses
a closure with its constant and linear terms removed:
\begin{equation}
  \mathcal N_\theta^{\circ}(\vect q)
  =\mathcal N_\theta(\vect q)-\mathcal N_\theta(\vect 0)
   -\left.\frac{\partial\mathcal N_\theta}{\partial\vect q}\right|_{\vect0}\vect q.
  \label{eq:origin-correction}
\end{equation}
Hence \(\mathcal N_\theta^{\circ}(\vect0)=\vect0\) and
\(\partial\mathcal N_\theta^{\circ}/\partial\vect q(\vect0)=\mat0\).
For concise notation, the superscript \(\circ\) is omitted below and
\(\mathcal N_\theta\) denotes the corrected closure.

The nonlinear PROM--ANN decoder is
\begin{equation}
  \widetilde{\vect u}(\vect q,\vect\mu)
  =\vect u_{\rm aff}(\vect\mu)
   +\mat V\mat A\vect q
   +\overline{\mat V}\mathcal N_\theta(\vect q).
  \label{eq:prom-ann-decoder}
\end{equation}
Its tangent with respect to the primary coordinates is the nonlinear-manifold
test basis
\begin{equation}
  \mat W(\vect q)
  =\frac{\partial\widetilde{\vect u}}{\partial\vect q}
  =\mat V\mat A
   +\overline{\mat V}
    \frac{\partial\mathcal N_\theta}{\partial\vect q}(\vect q)
  \in\R^{N_f\times3}.
  \label{eq:W}
\end{equation}
The restriction of \(\mat W\) to element \(e\) is denoted by \(\mat W_e\).

\section{Intrusive PROM--ANN and HPROM--ANN}
\label{sec:hprom}

\subsection{PROM--ANN Galerkin correction}

The intrusive reduced equilibrium equation is a Galerkin projection using the
tangent of the nonlinear decoder:
\begin{equation}
  \vect g(\vect q;\vect\mu)
  =\mat W(\vect q)\transpose
    \vect r\!\left(\widetilde{\vect u}(\vect q,\vect\mu);\vect\mu\right)
  =\vect0.
  \label{eq:prom-galerkin}
\end{equation}
The online solve initializes from \(\vect q^0(\vect\mu)\) in
\cref{eq:q0}, then applies reduced Newton iterations,
\begin{equation}
  \mat K_{\rm red}(\vect q^k)\Delta\vect q^k
  =-\vect g(\vect q^k;\vect\mu),
  \qquad
  \vect q^{k+1}=\vect q^k+\Delta\vect q^k.
  \label{eq:reduced-newton}
\end{equation}
Ignoring manifold curvature gives the standard Gauss--Newton-style reduced
tangent
\begin{equation}
  \mat K_{\rm GN}
  =\mat W\transpose\mat K\mat W.
  \label{eq:gn-tangent}
\end{equation}
The exact derivative of \cref{eq:prom-galerkin} also contains the curvature
term \(\mathrm D\mat W\transpose[\vect r]\).  The code has an option to
include that term.  This distinction is separate from MAW--ECM.

\subsection{Elementwise hyper-reduction}

Directly forming \cref{eq:prom-galerkin} still visits all \(N_e\) elements.
For a fixed empirical cubature rule with support \(\set Z_{\rm res}\) and
weights \(\xi_e^{\rm res}\), the hyper-reduced Galerkin residual is
\begin{equation}
  \vect g_{\rm HR}(\vect q;\vect\mu)
  \approx
  \sum_{e\in\set Z_{\rm res}}
  \xi_e^{\rm res}\,
  \mat W_e(\vect q)\transpose
  \vect r_e\!\left(\widetilde{\vect u}(\vect q,\vect\mu);\vect\mu\right).
  \label{eq:fixed-ecm-residual}
\end{equation}
Classical ECM selects a small positive support and weights offline, after a
randomized SVD compression of sampled elementwise operators.  Stage 9 builds
the residual data block for sampled state \(j\) as
\begin{equation}
  \mat Q_j^{\rm res}
  =\begin{bmatrix}
     \mat W_1\transpose\vect r_1 & \cdots &
     \mat W_{N_e}\transpose\vect r_{N_e}
   \end{bmatrix}
  \in\R^{3\times N_e},
  \qquad
  \vect b_j^{\rm res}=\mat Q_j^{\rm res}\vect1.
  \label{eq:ecm-residual-data}
\end{equation}
The column associated with element \(e\) is exactly its contribution to the
Galerkin residual.  Fixed ECM finds a support and weights that reproduce the
sampled blocks to the requested tolerance.

\subsection{D-HPROM--ANN}

The direct variant is not an equilibrium solve.  It uses
\begin{equation}
  \vect q=\vect q^0(\vect\mu),
  \qquad
  \widetilde{\vect u}_{\rm D}
  =\widetilde{\vect u}(\vect q^0(\vect\mu),\vect\mu),
  \label{eq:direct-state}
\end{equation}
then evaluates the homogenized quantities with the hyper-reduced output rule
in \cref{sec:hom-maw}.  It does not evaluate \cref{eq:fixed-ecm-residual} and
does not perform the correction \cref{eq:reduced-newton}.  Therefore its
quality relies on the affine initialization and ANN decoder being sufficiently
accurate over the deployment domain.

\section{MAW--ECM: state-dependent cubature weights}
\label{sec:maw}

\subsection{Core idea}

MAW--ECM keeps a \emph{common sparse element support} but allows its weights
to vary smoothly with the state.  For any target \(\tau\) (residual, strain,
or stress), let \(\set Z_\tau\) be its selected support and let
\(\vect\xi^\tau(\vect z)\) contain the support weights.  The generic
hyper-reduced operator is
\begin{equation}
  \sum_{e=1}^{N_e}\vect a_e^\tau
  \ \approx\
  \sum_{e\in\set Z_\tau}
    \xi_e^\tau(\vect z)\vect a_e^\tau.
  \label{eq:maw-generic}
\end{equation}
The support is fixed after offline pruning; only
\(\vect\xi^\tau(\vect z)\) changes online.  This is important for a reduced
mesh implementation because the mesh connectivity can stay fixed.

\subsection{Local MAW constraints and common-support pruning}

Start from a feasible classical ECM candidate support
\(\set Z_{\tau,\rm ini}\) with weights \(\vect\xi_{\rm ini}^\tau\).
At each training state \(\vect z_j\), build a local linear constraint
\begin{equation}
  \mat A_j^\tau\vect\xi_j^\tau=\vect b_j^\tau,
  \qquad
  \vect\xi_j^\tau\geq\vect0,
  \qquad
  \vect1\transpose\vect\xi_j^\tau=\alpha_\tau.
  \label{eq:maw-local-constraints}
\end{equation}
The last equality is used in the present final model with
\(\alpha_\tau=990\), the number of full RVE elements.  It preserves the
reference integration measure and makes the weights nonnegative and
comparable across states.

For residual MAW,
\begin{equation}
  \mat A_j^{\rm res}
  =\mat Q_j^{\rm res}(:,\set Z_{{\rm res},\rm ini}).
  \label{eq:maw-res-A}
\end{equation}
For strain and stress MAW, the elementwise homogenization data are
\begin{equation}
  \mat C_j^{\varepsilon}(:,e)=A_e\avg{\mat E_e}_e,
  \qquad
  \mat C_j^{\sigma}(:,e)=A_e\avg{\mat S_e}_e,
  \label{eq:maw-hom-data}
\end{equation}
and the corresponding \(\mat A_j^\tau\) is the restriction of the required
columns to the candidate support.  In a full-target variant,
\(\vect b_j^\tau\) is the sum over all elements.  In the final RVE MAW run,
\texttt{rhs\_mode=anchor}: the right hand side is
\begin{equation}
  \vect b_j^\tau
  =\mat A_j^\tau\vect\xi_{\rm ini}^\tau.
  \label{eq:maw-anchor-target}
\end{equation}
This preserves the classical ECM rule locally while MAW pruning seeks a
smaller common support and a smooth state-dependent weight field.

The pruning routine eliminates candidate elements while maintaining the
feasibility of \cref{eq:maw-local-constraints} over all sampled states.
Its optional second phase introduces state-space smoothness, schematically
\begin{equation}
  \min_{\{\vect\xi_j\}}\ 
  \sum_{(i,j)\in\set E_{\rm graph}}
  \norm{\vect\xi_i-\vect\xi_j}^2
  \quad\text{subject to the local constraints.}
  \label{eq:maw-smoothness}
\end{equation}
The exact pruning procedure is an active-set elimination algorithm; it
returns the common support and the feasible training matrix
\begin{equation}
  \mat\Xi_\tau
  =\begin{bmatrix}
    \vect\xi^\tau(\vect z_1)&\cdots&\vect\xi^\tau(\vect z_{N_\tau})
  \end{bmatrix}.
  \label{eq:maw-weight-training-data}
\end{equation}

\subsection{Weight regressors}

The support weights must be evaluated at new online states.  The final RVE
uses two regressor families.
\paragraph{Residual MAW--RBF.}
The residual support weights depend on the structured coordinates,
\(\vect z_{\rm res}=\vect q\).  A multi-output Gaussian RBF has the form
\begin{equation}
  \widehat{\vect\xi}^{\rm res}(\vect q)
  =\mat D\left[
     \mat\Phi(\vect q)\mat A_{\rm RBF}+\mat p(\vect q)\mat B_{\rm RBF}
    \right]\transpose,
  \qquad
  \Phi_\ell(\vect q)
  =\exp\!\left[-\frac12
      \sum_a\left(\frac{q_a-c_{\ell a}}{\ell_a}\right)^2\right].
  \label{eq:maw-rbf}
\end{equation}
Here \(\mat D\) restores output scaling and \(\mat p\) is an optional
constant or affine polynomial block.  The online implementation clips to
nonnegative values and renormalizes the weight sum to \(\alpha_{\rm res}\).
It also evaluates \(\partial\vect\xi^{\rm res}/\partial\vect q\)
analytically.

\paragraph{Homogenization MAW--ANN.}
The final strain and stress support weights depend on macro input,
\begin{equation}
  \vect z_{\varepsilon}=\vect z_{\sigma}=\vect\mu.
  \label{eq:hom-weight-input}
\end{equation}
Their ANN output is converted to positive, sum-preserving weights by
\begin{equation}
  \vect\xi^\tau(\vect\mu)
  =\alpha_\tau\,\operatorname{softmax}
    \left(\mathcal W^\tau_\omega(\vect\mu)\right),
  \qquad \tau\in\{\varepsilon,\sigma\}.
  \label{eq:maw-ann}
\end{equation}
This gives nonnegative weights with exactly the desired total by construction.

\paragraph{Current final MAW model.}
The saved file
\code{stage_12_hprom_ann_mawecm_res_eps_sig_phase1to40_phase2to10_sum990_ann/ecm_weights_all.npz}
uses a three-dimensional \(\vect q\), an RBF residual-weight map, ANN strain
and stress weight maps driven by \(\vect\mu\), and ten elements per target.

\section{HPROM--ANN with MAW--ECM}
\label{sec:maw-hprom}

\subsection{Dynamic residual and consistent tangent}

With state-dependent residual weights, the online projected residual becomes
\begin{equation}
  \vect g_{\rm MAW}(\vect q;\vect\mu)
  =\sum_{e\in\set Z_{\rm res}}
   \xi_e^{\rm res}(\vect q)
   \mat W_e(\vect q)\transpose
   \vect r_e\!\left(\widetilde{\vect u}(\vect q,\vect\mu);\vect\mu\right).
  \label{eq:maw-residual}
\end{equation}
Let \(\vect a_e=\mat W_e\transpose\vect r_e\).  Its exact directional
derivative has three contributions,
\begin{equation}
  \mathrm D\vect g_{\rm MAW}
  =\sum_{e\in\set Z_{\rm res}}
    \left[
      \xi_e^{\rm res}\,(\mathrm D\mat W_e\transpose)[\vect r_e]
      +\xi_e^{\rm res}\,\mat W_e\transpose\mat K_e\mat W_e
      +\vect a_e\left(\nabla_{\vect q}\xi_e^{\rm res}\right)\transpose
    \right].
  \label{eq:maw-tangent}
\end{equation}
The first term is the nonlinear-manifold curvature, the second is the usual
Galerkin stiffness term, and the third is the MAW weight-tangent term.  If the
weights are fixed, the last term vanishes.  If a Gauss--Newton manifold tangent
is used, the first term is neglected.  The production HPROM--ANN solver can
include both the manifold curvature and the MAW weight tangent; the
state-dependent weight tangent must not be silently omitted when using
\cref{eq:maw-residual}.

The Newton correction is
\begin{equation}
  \mat K_{\rm MAW}(\vect q^k)\Delta\vect q^k
  =-\vect g_{\rm MAW}(\vect q^k;\vect\mu),
  \qquad
  \vect q^{k+1}=\vect q^k+\Delta\vect q^k.
  \label{eq:maw-newton}
\end{equation}

\subsection{Hyper-reduced homogenization}
\label{sec:hom-maw}

After obtaining the accepted state, homogenized strain and stress are
approximated from their own selected supports:
\begin{align}
  \overline{\mat E}_{\rm HR}
  &\approx
  \frac{1}{A_0}
  \sum_{e\in\set Z_{\varepsilon}}
    \xi_e^{\varepsilon}(\vect\mu)
    A_e\avg{\mat E_e}_e,
  \label{eq:maw-hom-eps}\\
  \overline{\mat S}_{\rm HR}
  &\approx
  \frac{1}{A_0}
  \sum_{e\in\set Z_{\sigma}}
    \xi_e^{\sigma}(\vect\mu)
    A_e\avg{\mat S_e}_e.
  \label{eq:maw-hom-sig}
\end{align}
The residual, strain, and stress supports may overlap, but they have distinct
purposes and should be treated as independent rules.  A reduced HROM mesh is
formed from their union plus the nodes and conditions required by the strong
boundary conditions.

\subsection{Online algorithms}

\paragraph{HPROM--ANN with MAW--ECM.}
For each macro increment \(\vect\mu^n\):
\begin{enumerate}[label=\arabic*.,leftmargin=8mm]
  \item Compute \(\mat F(\vect\mu^n)\) and
  \(\vect u_{\rm aff}(\vect\mu^n)\).
  \item Form \(\vect q^0(\vect\mu^n)\) from \cref{eq:q0}.  Continuation
  from the previous converged \(\vect q\) may be used after the first step.
  \item Decode \(\widetilde{\vect u}\), evaluate \(\mat W\), and evaluate
  the residual MAW weights and their derivative.
  \item Assemble \cref{eq:maw-residual} and its selected tangent terms, then
  solve \cref{eq:maw-newton} until the reduced convergence criterion is met.
  \item Evaluate \cref{eq:maw-hom-eps}--\cref{eq:maw-hom-sig} with the
  accepted state and the strain/stress MAW weights.
\end{enumerate}

\paragraph{D-HPROM--ANN.}
For each macro increment \(\vect\mu^n\):
\begin{enumerate}[label=\arabic*.,leftmargin=8mm]
  \item Compute \(\vect q=\vect q^0(\vect\mu^n)\).
  \item Decode \(\widetilde{\vect u}_{\rm D}\) through
  \cref{eq:prom-ann-decoder}.
  \item Evaluate only the hyper-reduced homogenization
  \cref{eq:maw-hom-eps}--\cref{eq:maw-hom-sig}.
\end{enumerate}
This direct path is appropriate only after verifying that its lack of
equilibrium correction is acceptable on unseen trajectories.

\section{Offline workflow and source map}
\label{sec:source-map}

\begin{longtable}{@{}p{16mm}p{48mm}p{90mm}@{}}
\toprule
Stage & Main files & Role \\
\midrule
\endhead
0 & \code{stage0_training_trajectory.py} & Builds the ten serpentine macro-strain paths and writes \code{stage_0_trajectories.npz}. \\
1 & \code{stage1_training_fom_solver_rve.py}, \code{fom_solver_rve.py} & Runs the FOM for each path and saves displacement, macro-strain, homogenized strain, and homogenized stress histories. \\
2 & \code{stage2_pod_rve.py} & Removes \(\vect u_{\rm aff}\), checks strong Dirichlet consistency, and computes \(\mat V_{\rm tot}\). \\
7a & \code{stage7a_prepare_rbf_dataset_ls.py}; wrapper \code{stage7a_prepare_ann_dataset_ls.py} & Implements \cref{eq:fit-T}--\cref{eq:Baff}; saves \(\mat V\), \(\overline{\mat V}\), \(\mat A\), \(\mat B_{\rm aff}\), and ANN data. \\
7b & \code{stage7b_train_ann_manifold.py} & Trains \(\mathcal N_\theta:\vect q\mapsto\overline{\vect q}\). \\
9a & \code{stage9_build_ecm_dataset_ann.py} & Builds sampled elementwise residual and homogenization operators. \\
9b & \code{stage9_compute_ecm_weights_ann.py} & Builds the classical fixed ECM rule used as the MAW anchor. \\
10 & \code{stage10_test_hprom_ann_ls.py}, \code{hprom_ann_solver_rve.py} & Runs and benchmarks PROM--ANN, HPROM--ANN, and D-HPROM--ANN. \\
12 & \code{stage12b_build_mawecm_model_gpr.py}, \code{mawecm_pruning.py}, \code{mawecm_rbf_weights.py}, \code{mawecm_ann_weights.py} & Builds the common supports, local feasible MAW weights, and online RBF/ANN weight regressors.  The builder name is historical; it supports the ANN workflow. \\
\bottomrule
\end{longtable}

The two most important persisted artifacts for the current case are:
\begin{itemize}[leftmargin=7mm]
  \item \code{stage_7_ann_data_ls/}: structured coordinates,
  primary/secondary bases, \(\mat A\), \(\mat T\), and the affine
  initialization fit;
  \item \code{stage_12_hprom_ann_mawecm_res_eps_sig_phase1to40_phase2to10_sum990_ann/ecm_weights_all.npz}:
  all MAW supports, models, HROM mesh metadata, and diagnostic data.
\end{itemize}

\section{Validation logic and error decomposition}
\label{sec:validation}

The method has several distinct approximation layers.  They should be
validated separately before interpreting a final stress error.
\begin{enumerate}[label=\arabic*.,leftmargin=8mm]
  \item \textbf{POD truncation:} compare \(\vect s\) with
  \(\mat V_{\rm tot}\mat V_{\rm tot}\transpose\vect s\).
  \item \textbf{Coordinate construction:} check
  \(\frobnorm{\mat T\mat Q_{\rm tot}-\mat M}\), the invertibility and
  conditioning of \(\mat A\), and the exact split reconstruction in
  \cref{eq:split-reconstruction}.
  \item \textbf{Affine initialization:} evaluate the unseen-trajectory error
  of \(\vect q^0(\vect\mu)\) from \cref{eq:q0}.
  \item \textbf{ANN closure:} evaluate the error in
  \(\overline{\vect q}\), in the decoded state, and in the nonlinear tangent
  \(\mat W\) if it is used intrusively.
  \item \textbf{Fixed ECM and MAW constraints:} verify every sampled local
  constraint \(\mat A_j\vect\xi_j=\vect b_j\), positivity, and the required
  weight sum.  Validate regressed weights away from the fitting nodes.
  \item \textbf{HPROM correction:} compare the corrected \(\vect q\) and
  reduced residual convergence against the non-hyper-reduced PROM--ANN.
  \item \textbf{Output accuracy:} compare FOM, PROM--ANN, HPROM--ANN, and
  D-HPROM--ANN homogenized stress histories on trajectories excluded from
  the offline construction.
\end{enumerate}

This order localizes failures.  For example, a good ANN state with a poor
stress curve points to output hyper-reduction, whereas a poor direct state but
good corrected HPROM state points to \(\vect q^0\), not to MAW--ECM.

\section{Porting the technology to Burgers 1D}
\label{sec:burgers}

\subsection{What transfers unchanged}

The reusable architecture is independent of the RVE:
\begin{equation}
  \begin{gathered}
    \text{training trajectories}
    \longrightarrow \text{snapshots}
    \longrightarrow \text{POD}
    \longrightarrow \text{structured coordinates}
    \\[0.3em]
    \longrightarrow \text{ANN closure}
    \longrightarrow \text{Galerkin correction and MAW--ECM}.
  \end{gathered}
  \label{eq:portable-chain}
\end{equation}
For Burgers, replace the RVE state, residual, and homogenized stress by the
finite-element or finite-volume Burgers state, time-discrete residual, and
the desired output functional.  The coordinate construction,
\(\mat V\mat A\vect q+\overline{\mat V}\mathcal N_\theta(\vect q)\), and
the elementwise MAW residual construction remain the same.

\subsection{Transient residual}

For example, a backward-Euler semi-discretization of Burgers at time step
\(n\) has a residual of the generic form
\begin{equation}
  \vect r^n(\vect u^n;\vect u^{n-1},\vect\mu^n)
  =\frac{1}{\Delta t}\mat M_h(\vect u^n-\vect u^{n-1})
   +\vect f_{\rm conv}(\vect u^n)
   +\vect f_{\rm diff}(\vect u^n;\vect\mu^n)
   -\vect f_{\rm ext}(\vect\mu^n).
  \label{eq:burgers-residual}
\end{equation}
The PROM--ANN Galerkin residual becomes
\begin{equation}
  \vect g^n(\vect q^n)
  =\mat W(\vect q^n)\transpose
    \vect r^n\!\left(
      \widetilde{\vect u}(\vect q^n,\vect\mu^n);
      \widetilde{\vect u}(\vect q^{n-1},\vect\mu^{n-1}),
      \vect\mu^n
    \right).
  \label{eq:burgers-prom}
\end{equation}
The elementwise residual operator used by ECM must include both the spatial
terms and the time-discrete mass contribution assigned to each element.

\subsection{Critical modeling decision for a transient problem}

For the RVE, \(\vect\mu\) identifies a quasi-static state well enough that
the affine map \(\vect q^0(\vect\mu)\) is highly accurate.  This is not
automatically true for Burgers.  Two states can have the same viscosity or
forcing parameter but differ because of time history.  Before porting
\cref{eq:Baff}, decide what identifies the current state:
\begin{enumerate}[label=\arabic*.,leftmargin=8mm]
  \item If the dynamics are deterministic from a known initial condition,
  include time and the active forcing descriptors in \(\vect\mu^n\), then fit
  \(\vect q^0(\vect\mu^n)\).
  \item If history matters beyond the chosen input, use continuation
  \(\vect q^{0,n}=\vect q^{n-1}\) as the principal initializer and retain the
  intrusive reduced solve.
  \item If a direct model is needed, augment the ANN input with enough state
  history or learn a recurrent/time-step map.  Do not claim a direct map from
  an insufficient parameter vector to a history-dependent state.
\end{enumerate}

\subsection{Burgers-specific offline plan}

\begin{enumerate}[label=\arabic*.,leftmargin=8mm]
  \item Define parameter/time trajectories that cover viscosity, forcing,
  boundary conditions, initial-condition families, and the time interval of
  interest.
  \item Generate FOM snapshots.  If inhomogeneous Dirichlet conditions are
  present, remove an exact lifting as in \cref{eq:affine-lifting}; otherwise
  use the state directly.
  \item Compute a POD basis and select the primary dimension from the
  parameter/state dimension that is physically meaningful for the intended
  latent representation.
  \item Fit \(\mat T\) using \cref{eq:fit-T}, factor
  \(\mat V_{\rm tot}\mat T\transpose=\mat V\mat A\), form \(\vect q\), and
  train the secondary-coordinate closure.
  \item Build residual ECM blocks using
  \(\mat W_e\transpose\vect r_e^n\) at sampled time states.  Build a separate
  output ECM block only if a spatial integral, flux, energy, or other output
  requires it.
  \item Build classical ECM first, then MAW supports and weight regressors.
  The residual MAW input should normally be \(\vect q\); an output MAW input
  can be \(\vect q\), \(\vect\mu\), or both, provided it identifies the
  integrand reliably.
  \item Validate fixed PROM, fixed HPROM, MAW HPROM, and a direct variant on
  unseen trajectories, including long-time stability and residual convergence.
\end{enumerate}

\section{Implementation checklist for a future coding session}
\label{sec:checklist}

Before implementing the Burgers version, establish the following explicitly.
\begin{enumerate}[label=\arabic*.,leftmargin=8mm]
  \item Snapshot orientation: this document uses states as columns in
  \(\mat S\), whereas several NumPy files store samples as rows.  Transpose
  equations consistently rather than mixing conventions.
  \item Lifting: define whether the reduced state is the full field or an
  affine-free fluctuation.  Do not build a basis from mixed boundary data.
  \item Coordinate definition: retain the exact sequence
  \(\mat Q_{\rm tot}\), \(\mat T\), \(\mat V\mat A\), then \(\mat Q\).
  Avoid introducing a second, ambiguously named ``master'' coordinate.
  \item Decoder: use
  \(\widetilde{\vect u}=\vect u_{\rm aff}+\mat V\mat A\vect q+
  \overline{\mat V}\mathcal N_\theta(\vect q)\), and differentiate it to
  obtain the Galerkin test basis \(\mat W\).
  \item Direct versus intrusive path: the direct path is an approximation;
  the intrusive path enforces a projected residual.  Benchmark both instead
  of assuming that a good closure automatically gives a good direct solver.
  \item MAW: supports are fixed, weights are dynamic.  Include
  \(\partial\vect\xi/\partial\vect q\) in the reduced tangent when residual
  weights depend on \(\vect q\).
  \item Outputs: use a separate hyper-reduced rule for every quantity whose
  elementwise integrand differs materially from the residual.
\end{enumerate}

\section{Compact formula reference}
\label{sec:formula-reference}

\begin{align}
  \mat E &= \tfrac12(\mat F\transpose\mat F-\mat I),
  & \vect\mu&=[E_{xx},E_{yy},G_{xy}]\transpose,
  \\
  \vect u&=\vect u_{\rm aff}+\vect s,
  & \mat Q_{\rm tot}&=\mat V_{\rm tot}\transpose\mat S,
  \\
  \mat T&=\operatorname*{arg\,min}_{\widehat{\mat T}}
  \frobnorm{\widehat{\mat T}\mat Q_{\rm tot}-\mat M}^{2},
  & \mat V_{\rm tot}\mat T\transpose&=\mat V\mat A,
  \\
  \mat Q&=\mat A^{-1}\mat V\transpose\mat S,
  & \overline{\mat Q}&=\overline{\mat V}\transpose\mat S,
  \\
  \mat B_{\rm aff}&=\operatorname*{arg\,min}_{\mat B}
  \frobnorm{[\mat M\transpose,\vect1]\mat B-\mat Q\transpose}^{2},
  & (\vect q^0)\transpose&=[\vect\mu\transpose,1]\mat B_{\rm aff},
  \\
  \widetilde{\vect u}(\vect q,\vect\mu)
  &=\vect u_{\rm aff}(\vect\mu)+\mat V\mat A\vect q+
  \overline{\mat V}\mathcal N_\theta(\vect q),
  & \mat W&=\frac{\partial\widetilde{\vect u}}{\partial\vect q},
  \\
  \vect g_{\rm MAW}
  &=\sum_{e\in\set Z_{\rm res}}
  \xi_e^{\rm res}(\vect q)\mat W_e\transpose\vect r_e,
  & \mat K_{\rm MAW}\Delta\vect q&=-\vect g_{\rm MAW},
  \\
  \overline{\mat S}_{\rm HR}
  &\approx\frac{1}{A_0}\sum_{e\in\set Z_\sigma}
   \xi_e^\sigma(\vect\mu)A_e\avg{\mat S_e}_e.
\end{align}

\end{document}
